\section{Tropical feasibility and fixed corank}
\label{sec:tropical}

Akian, Gaubert, and Guterman identify the coordinates that can be finite
in a solution of tropical inequalities with the winning states of an
associated mean-payoff game
\cite[Theorem~3.2 and Corollary~3.7]{AkianGaubertGuterman2012}.
Together with Theorem~\ref{thm:main}, this gives binary-length bounds
for tropical feasibility and for rank within a fixed distance of the
number of columns. We use the conventions $a+(-\infty)=-\infty$
and $\max\varnothing=-\infty$.

\begin{corollary}[Tropical feasibility and fixed corank]
\label{cor:tropical-computation}
For explicitly listed $m\times n$ matrices $A,B$ and $m$-dimensional
vectors $c,d$ over $\mathbb Q\cup\{-\infty\}$, feasibility of
\[
 \max\bigl\{c_i,\max_j(A_{ij}+x_j)\bigr\}
 \leq \max\bigl\{d_i,\max_j(B_{ij}+x_j)\bigr\}
 \qquad\text{for every row }i
\]
is decidable deterministically in
$2^{O((\log(L+2))^2)}$ bit operations, where $L$ is the complete
explicit binary input length. Variables range over
$\mathbb R\cup\{-\infty\}$; feasibility with all coordinates finite
is decidable within the same bound. Equalities between the two displayed
maxima are allowed, interpreted as pairs of inequalities.
For a homogeneous system ($c=d=-\infty$), one can also compute the
maximal finite-coordinate support of a solution, including the empty
support. In particular, the same bound decides feasibility of max-atom
systems $z\leq\max(x,y)+a$, with rational constants $a$ and finite
real, equivalently rational, variables.

For each fixed integer $k\geq0$, a deterministic algorithm with bit
bound $2^{O_k((\log(L+2))^2)}$ decides whether an explicitly listed
$m\times n$ matrix over $\mathbb Q\cup\{-\infty\}$, with $m\geq n\geq1$,
has tropical rank at least $n-k$. Here tropical rank is the largest
order of a square submatrix whose permutation sums have a finite,
unique maximum.
The rank is zero if no such nonempty submatrix exists.
\end{corollary}

\begin{proof}
\noindent\textbf{Integer data and homogeneous supports.}\enspace
Let $Q$ be the product of the positive denominators of the finite input
coefficients, taking $Q=1$ when there are none. Multiplying every finite
coefficient and variable by $Q$ preserves the inequalities and the sets
of finite coordinates. It also preserves the comparisons between
permutation sums that define tropical rank. The bit length of $Q$ is
at most one plus the sum of the denominator bit lengths, so the scaled integer
data have polynomial binary length.

Consider first a homogeneous system
\begin{equation}
 \max_j(C_{ij}+z_j)\leq\max_j(D_{ij}+z_j)
 \qquad\text{for every }i.
 \label{eq:tropical-homogeneous}
\end{equation}
The support of $z$ is $\{j:z_j\ne-\infty\}$. To obtain a total game,
apply the preprocessing of \cite[Section~2.1]{AkianGaubertGuterman2012}.
Discard a row whose left side is identically $-\infty$. If a right
side is identically $-\infty$, every variable with a finite coefficient
on that row's left side must equal $-\infty$. Record these coordinates,
substitute $-\infty$ for them throughout the system, and repeat.
Each step removes a row or a variable, so this takes polynomial time.
If no variables remain, the maximal support is empty. Otherwise add
the tautologies $z_j\leq z_j$ for the remaining variables.

The resulting bipartite game has a Min vertex for each remaining column
and a Max vertex for each row. A finite $C_{ij}$ gives an edge from
column $j$ to row $i$ of weight $-C_{ij}$; a finite $D_{ij}$ gives
an edge from row $i$ to column $j$ of weight $D_{ij}$. Every vertex
has an outgoing edge. Theorem~3.2 of
\cite{AkianGaubertGuterman2012} identifies its winning column vertices
with the maximal support of a solution. The recorded deleted coordinates
belong to no support. The union of all supports is attained, because
coordinatewise maxima of finitely many solutions are again solutions;
the identically $-\infty$ solution realizes the empty case.

The cited game counts a column--row--column pair as one turn. Counting
the two edges separately halves the mean payoff; the contribution of
an unmatched edge tends to zero. The source uses upper limiting averages;
positional optimality makes their nonnegative winning set coincide with
that for our lower limiting averages. The game has
polynomially many vertices and edges, with weights copied or negated
from the scaled coefficients or equal to zero. Its full binary encoding
is polynomial in $L+1$, so Theorem~\ref{thm:main} computes the maximal
support within the asserted bound.

\par\medskip\noindent\textbf{Affine and finite feasibility.}\enspace
After scaling, append $c$ to $A$ and $d$ to $B$ as a new column,
and introduce one coordinate $t$. The homogeneous system is
\[
 \max\bigl\{c_i+t,\max_j(A_{ij}+z_j)\bigr\}
 \leq\max\bigl\{d_i+t,\max_j(B_{ij}+z_j)\bigr\}.
\]
Every affine solution $x$ gives a homogeneous solution $(x,0)$.
Conversely, a homogeneous solution with finite $t$ gives the affine
solution $x_j=z_j-t$. Thus affine feasibility asks whether the new
coordinate belongs to the maximal support. Feasibility with all original
coordinates finite asks whether that support contains all $n+1$
coordinates. A coordinate eliminated during preprocessing is absent
from the computed support, so any test requiring it fails; this includes
$t$ if it is forced to $-\infty$.
This is the homogenization underlying
\cite[Theorem~3.5 and Corollaries~3.4 and~3.7]{AkianGaubertGuterman2012}.
Equalities are handled by imposing both inequalities.

A max-atom $z\leq\max(x,y)+a$ is a homogeneous row with one term on
the left and two on the right. Apply the all-finite test. With integer
coefficients, taking the floor of each finite coordinate preserves every
inequality: flooring commutes with integer translation and with finite
maxima. Scaling back therefore gives a rational solution whenever a real
one exists. This also explains the integral-solution interface in
\cite[Proposition~3.9, Corollary~3.10 and Remark~3.15]
{AkianGaubertGuterman2012}.

\par\medskip\noindent\textbf{Fixed corank.}\enspace
If $n\leq k$, the rank threshold is nonpositive and the answer is yes.
Otherwise set $r=n-k$. The rank criterion of
\cite[Corollaries~4.16--4.17]{AkianGaubertGuterman2012} says that the
rank is at least $r$ exactly when some set of $r$ columns is tropically
independent. For each such set let $M$ be the resulting $m\times r$
matrix. Its columns are dependent precisely when there is a vector
$z\not\equiv-\infty$ such that every row maximum is attained at least
twice or equals $-\infty$. Equivalently, $z$ satisfies the homogeneous
system
\begin{equation}
 M_{ij}+z_j\leq\max_{\ell\ne j}(M_{i\ell}+z_\ell)
 \qquad(1\leq i\leq m,\ 1\leq j\leq r).
 \label{eq:tropical-dependence}
\end{equation}
Its maximal support is therefore empty exactly when the selected columns
are independent. This is the inequality form of the game reduction in
\cite[Theorem~4.9]{AkianGaubertGuterman2012}. The row index $(i,j)$
retains the column excluded from Max's next choice. Empty right sides
and all-$-\infty$ columns are covered by the support computation above,
including when $r=1$.

There are at most $mr$ inequalities, with $r$ coefficients on each
side, so each test has polynomial encoding length. There are
$\binom nr=\binom nk\leq n^k$ tests. For fixed $k$, their total
bit cost is $2^{O_k((\log(L+2))^2)}$, as claimed.
\end{proof}

Strongly polynomial algorithms have been claimed using max-atoms and
tropical optimization
\cite{Truffet2025MaxAtoms,Truffet2026SubstitutionV4}.
Appendix~A of the randomized companion gives a two-variable
counterexample to the printed variable-selection rule in
\texttt{arXiv:2603.26423v4}~\cite{OpenAI2026Randomized}.
That comparison is specific to the version and rule examined; it
neither assesses the distinct 2025 max-atom proof nor transfers to
a revised procedure. It has no role in the proofs of this article.
